% Lesion-level joint conformal risk control — MICCAI-style (LNCS) submission draft.
% Compile: latexmk -pdf main.tex   (or upload the paper/ folder to Overleaf)
% All numbers come from \newcommand{\FracatlasUnits}{719}
\newcommand{\FracatlasMulti}{159}
\newcommand{\FracatlasWithSmall}{234}
\newcommand{\FracatlasSmallCut}{104}
\newcommand{\FracatlasSeeds}{3}
\newcommand{\FracatlasReps}{1000}
\newcommand{\FracatlasDice}{0.057}
\newcommand{\FracatlasHOne}{INCONCLUSIVE}
\newcommand{\FracatlasHTwoFive}{FAILED}
\newcommand{\FracatlasHTwoOnezero}{FAILED}
\newcommand{\FracatlasHTwoOnefive}{FAILED}
\newcommand{\FracatlasLttThresholdRel}{+21.5\%}
\newcommand{\FracatlasOursUnionRel}{+1.8\%}
\newcommand{\FracatlasPixelCrcRel}{-0.2\%}
\newcommand{\FracatlasUnionDilationRel}{+16.4\%}
\newcommand{\KvasirSegUnits}{1000}
\newcommand{\KvasirSegMulti}{48}
\newcommand{\KvasirSegWithSmall}{312}
\newcommand{\KvasirSegSmallCut}{3774}
\newcommand{\KvasirSegSeeds}{3}
\newcommand{\KvasirSegReps}{1000}
\newcommand{\KvasirSegDice}{0.413}
\newcommand{\KvasirSegHOne}{PROCEED}
\newcommand{\KvasirSegHTwoFive}{FAILED}
\newcommand{\KvasirSegHTwoOnezero}{FAILED}
\newcommand{\KvasirSegHTwoOnefive}{FAILED}
\newcommand{\KvasirSegLttThresholdRel}{+10.2\%}
\newcommand{\KvasirSegOursUnionRel}{+0.7\%}
\newcommand{\KvasirSegPixelCrcRel}{-0.1\%}
\newcommand{\KvasirSegUnionDilationRel}{+1.6\%}
\newcommand{\SimVerdict}{FAILED}
\newcommand{\SimBest}{-47.2\%}
\newcommand{\SimBestGamma}{0.5}
\newcommand{\SimWorst}{+4.0\%}
\newcommand{\ResultValidity}{Every method that claims a guarantee met it on held-out data}
\newcommand{\ResultAbstract}{Every method that claims a guarantee met it on held-out data. Joint lesion-level control cost at most 0.7\% in set size over pixel-only CRC. Learn-Then-Test (a high-probability guarantee) needed 10--46\% larger sets and morphological sets 1--43\% larger; CCRA-S strata ranged -2 to +16\%. The pre-registered claim that lesion-equalized sets are smaller than joint thresholding was not supported on FracAtlas; not supported on Kvasir-SEG. In a controlled simulation they reduced set size by up to 47\% when small lesions are moderately less confident, and fell back to thresholding when lesions are indistinguishable from background.}
 (written by
% backend/scripts/experiments/paper_macros.py from docs/research/results/).
\documentclass[runningheads]{llncs}
\usepackage[T1]{fontenc}
\usepackage{amsmath,amssymb}
\usepackage{graphicx,booktabs,multirow,xcolor}
\usepackage{pgfplots,pgfplotstable}
\pgfplotsset{compat=1.17}
\usepackage[hidelinks]{hyperref}

\newcommand{\FracatlasUnits}{719}
\newcommand{\FracatlasMulti}{159}
\newcommand{\FracatlasWithSmall}{234}
\newcommand{\FracatlasSmallCut}{104}
\newcommand{\FracatlasSeeds}{3}
\newcommand{\FracatlasReps}{1000}
\newcommand{\FracatlasDice}{0.057}
\newcommand{\FracatlasHOne}{INCONCLUSIVE}
\newcommand{\FracatlasHTwoFive}{FAILED}
\newcommand{\FracatlasHTwoOnezero}{FAILED}
\newcommand{\FracatlasHTwoOnefive}{FAILED}
\newcommand{\FracatlasLttThresholdRel}{+21.5\%}
\newcommand{\FracatlasOursUnionRel}{+1.8\%}
\newcommand{\FracatlasPixelCrcRel}{-0.2\%}
\newcommand{\FracatlasUnionDilationRel}{+16.4\%}
\newcommand{\KvasirSegUnits}{1000}
\newcommand{\KvasirSegMulti}{48}
\newcommand{\KvasirSegWithSmall}{312}
\newcommand{\KvasirSegSmallCut}{3774}
\newcommand{\KvasirSegSeeds}{3}
\newcommand{\KvasirSegReps}{1000}
\newcommand{\KvasirSegDice}{0.413}
\newcommand{\KvasirSegHOne}{PROCEED}
\newcommand{\KvasirSegHTwoFive}{FAILED}
\newcommand{\KvasirSegHTwoOnezero}{FAILED}
\newcommand{\KvasirSegHTwoOnefive}{FAILED}
\newcommand{\KvasirSegLttThresholdRel}{+10.2\%}
\newcommand{\KvasirSegOursUnionRel}{+0.7\%}
\newcommand{\KvasirSegPixelCrcRel}{-0.1\%}
\newcommand{\KvasirSegUnionDilationRel}{+1.6\%}
\newcommand{\SimVerdict}{FAILED}
\newcommand{\SimBest}{-47.2\%}
\newcommand{\SimBestGamma}{0.5}
\newcommand{\SimWorst}{+4.0\%}
\newcommand{\ResultValidity}{Every method that claims a guarantee met it on held-out data}
\newcommand{\ResultAbstract}{Every method that claims a guarantee met it on held-out data. Joint lesion-level control cost at most 0.7\% in set size over pixel-only CRC. Learn-Then-Test (a high-probability guarantee) needed 10--46\% larger sets and morphological sets 1--43\% larger; CCRA-S strata ranged -2 to +16\%. The pre-registered claim that lesion-equalized sets are smaller than joint thresholding was not supported on FracAtlas; not supported on Kvasir-SEG. In a controlled simulation they reduced set size by up to 47\% when small lesions are moderately less confident, and fell back to thresholding when lesions are indistinguishable from background.}


\begin{document}
\title{Pixel Guarantees Hide Missed Lesions:\\Lesion-Level Joint Conformal Risk Control\\with Lesion-Equalized Prediction Sets}
\titlerunning{Lesion-Level Joint Conformal Risk Control}
\author{Mahendra Sajjan Choudhary \and Vansh Bhatiya \and Shashank Mouli Satapathy}
\authorrunning{M. S. Choudhary et al.}
\institute{School of Computer Science and Engineering, Vellore Institute of Technology, Vellore, India}
\maketitle

\begin{abstract}
Conformal risk control gives segmentation models a finite-sample bound on the
pixel false-negative rate. That bound is carried by large lesions and says
nothing about whether small lesions are found. We define lesion-miss and
small-lesion-miss risks on ground-truth connected components, prove they are
monotone under any nested prediction-set family, and control them
simultaneously with pixel FNR in expectation, with no Bonferroni split. We then
show that thresholding is an inefficient family for a lesion risk, because the
least confident lesion sets the global cutoff. We propose
\emph{lesion-equalized} prediction sets, which peak-normalise every candidate
lesion. The family is amplitude-invariant and excludes background above the
candidate level, and choosing it on a tuning split preserves the guarantee
exactly. In a pre-registered evaluation on FracAtlas and Kvasir-SEG with six
baselines, and in a controlled simulation: \ResultAbstract
\keywords{Conformal risk control \and Segmentation \and Lesion detection \and Uncertainty quantification}
\end{abstract}

\section{Introduction}

Conformal risk control (CRC)~\cite{angelopoulos2024crc} has become the standard
way to give a segmentation model a finite-sample guarantee. A threshold is
calibrated on held-out data so that the expected false-negative rate (FNR) of
the resulting prediction set stays below a user-chosen $\alpha$. Recent work
adapts this to 3-D lesion segmentation~\cite{tan2025cls}, to
image-conditional risk~\cite{luo2025ccra,recirc2026}, to morphological
sets~\cite{mossina2025morph} and to per-organ recall~\cite{perorgan2026}.

All of these guarantees are stated in \emph{pixels} (or voxels). Clinically,
the unit of harm is usually the \emph{lesion}. A missed 4\,mm metastasis next to
a well-segmented 4\,cm primary costs almost nothing in pixel FNR, yet it is
exactly the error a radiologist cares about. A pixel guarantee is carried by
large lesions and says nothing about small ones. Even methods motivated by small
lesions~\cite{tan2025cls} formalise their risk as a voxel fraction.

We make three contributions.
\begin{enumerate}
\item \textbf{Lesion-level joint risk control.} We define lesion-miss and
small-lesion-miss losses on ground-truth connected components. We show they are
monotone under every nested set family, and control them \emph{simultaneously}
with pixel FNR using a union or max-form CRC threshold. This requires neither a
Bonferroni split nor a high-probability bound (\S\ref{sec:method}).
\item \textbf{Lesion-equalized prediction sets.} We show why thresholding is
inefficient for a lesion loss and propose a nested family that peak-normalises
each candidate lesion. It is amplitude-invariant, it excludes background above
the candidate level, and it contains thresholding as a special case. The family
is chosen on a tuning split, which keeps the guarantee exact.
\item \textbf{A pre-registered evaluation.} We measure the pixel-to-lesion gap
and compare against six guarantee-matched or guarantee-weaker baselines (pixel
CRC, joint thresholding, max-form, Learn-Then-Test~\cite{angelopoulos2021ltt},
morphological sets~\cite{mossina2025morph} and CCRA-S strata~\cite{luo2025ccra})
on two public datasets and a controlled simulation. The protocol and decision
rules were committed before any result existed, and negative outcomes are
reported as such.
\end{enumerate}

\section{Related work}

\paragraph{Risk control for segmentation.}
CRC~\cite{angelopoulos2024crc} bounds the expectation of any monotone bounded
loss, and segmentation examples use pixel FNR. CLS~\cite{tan2025cls} calibrates
a per-volume voxel-FNR tolerance with a high-probability guarantee. Despite its
lesion-level motivation, its loss is a voxel fraction, so a missed small lesion
weighs in proportion to its volume. CRA/CCRA/CCRA-S~\cite{luo2025ccra} and
ReCIRC~\cite{recirc2026} target image-conditional pixel risk through adaptive
scores, strata or per-input budgets. Per-organ recall control~\cite{perorgan2026}
calibrates voxel FNR per organ, with Bonferroni across organs. FDR control for
instance segmentation~\cite{fdr2025} controls pixel FDR or FNR, one at a time.
None of these defines the loss on ground-truth components or controls several
granularities jointly.

\paragraph{Structured prediction sets.}
Morphological prediction sets~\cite{mossina2025morph} replace thresholding with
dilations of the predicted mask, and a follow-up controls false positives by
erosion~\cite{mossina2025fp}. Both change the nested family to suit
\emph{containment}. We change it to suit a \emph{lesion} loss, and we include
dilation as a baseline.

\paragraph{Multiple risks.}
Learn-Then-Test~\cite{angelopoulos2021ltt} controls many risks with
family-wise error control and needs no monotonicity, at the cost of a
high-probability guarantee and more conservative thresholds. Andéol et
al.~\cite{andeol2023} use box-wise recall for detection, one risk at a time.
Our joint guarantee is the special case where all losses are monotone and only
expectation control is required. It is simpler and tighter, and we measure the
difference.

\section{Method}\label{sec:method}

\paragraph{Setup.}
A segmenter outputs a posterior $\hat p_i:\Omega\to[0,1]$ for unit $i$ (an image
or a volume). The ground truth $Y_i\subseteq\Omega$ has connected components
$Y_i^1,\dots,Y_i^{K_i}$, which are its \emph{lesions}. A \emph{family} is a fixed,
label-free score transform $s_i=T(\hat p_i)\in[0,1]^\Omega$ with prediction sets
$S_\lambda(i)=\{u: s_i(u)\ge 1-\lambda\}$, $\lambda\in[0,1]$. These sets are nested
in $\lambda$ by construction.

\paragraph{Losses.}
For a hit fraction $\tau\in(0,1]$,
\begin{align*}
L^{\mathrm{pix}}_i(\lambda) &= 1-\tfrac{|S_\lambda\cap Y_i|}{|Y_i|}, &
L^{\mathrm{les}}_i(\lambda) &= \tfrac{1}{K_i}\textstyle\sum_k \mathbb 1\!\left[|S_\lambda\cap Y_i^k|<\tau|Y_i^k|\right],
\end{align*}
and $L^{\mathrm{small}}_i$ is $L^{\mathrm{les}}_i$ restricted to lesions below a
size cut fixed from the label distribution. Units without a lesion (or without a
small lesion) are excluded from the corresponding loss, so each guarantee is
conditional on the structure being present. All three losses are bounded by $1$
and are of the form $f(S_\lambda\cap Y_i)$ with $f$ non-increasing under
inclusion. They are therefore non-increasing in $\lambda$ for \emph{every}
nested family (Lemma~1, App.~\ref{app:proofs}). The lesion loss reduces to one
order statistic per lesion: lesion $k$ is hit at cutoff $c$ iff the
$\lceil\tau|Y^k|\rceil$-th largest score in $Y^k$ is $\ge c$ (Lemma~2). This makes
it as cheap to evaluate as the pixel loss.

\paragraph{Joint guarantee.}
Let $\hat\lambda^r$ be the split-CRC threshold~\cite{angelopoulos2024crc} for
risk $r\in\{\mathrm{pix},\mathrm{les},\mathrm{small}\}$, i.e.\ the smallest
$\lambda$ with $(n_r\hat R^r_{n_r}(\lambda)+1)/(n_r+1)\le\alpha$. The union
threshold $\hat\lambda^\cup=\max_r\hat\lambda^r$ satisfies
$\mathbb E[L^r_{n+1}(\hat\lambda^\cup)]\le\alpha$ for every $r$ simultaneously.
By monotonicity it is at least as large as each individual threshold, and no
Bonferroni split is needed because each bound is an expectation bound controlled
on its own. Applying CRC to $\max_r L^r_i$ instead bounds the expected
\emph{per-unit worst} risk. This is strictly stronger, at a larger $\lambda$.

\paragraph{Lesion-equalized prediction sets.}
Thresholding ($s=\hat p$) is the natural family for the pixel loss, but it is the
wrong one for a lesion loss. Under thresholding, lesion $k$ is hit only if a
$\tau$-fraction of its posterior exceeds $c$. The lesion constraint therefore
pushes the global cutoff below the \emph{least confident} lesions, and those are
disproportionately small. Every background pixel above that low cutoff then
enters every set. We instead peak-normalise each \emph{candidate} lesion. Let
$C_1,\dots,C_J$ be the connected components of $\{\hat p\ge t_0\}$ with peaks
$P_j=\max_{C_j}\hat p$. Then
\begin{equation}\label{eq:le}
s_{\rho,t_0}(u)=\begin{cases}\hat p(u)\,P_j^{-\rho} & u\in C_j,\\ \hat p(u) & \text{otherwise,}\end{cases}\qquad \rho\in[0,1].
\end{equation}
Three properties drive the method (proofs in App.~\ref{app:proofs}).
(i)~$s\in[0,1]$, every candidate pixel outranks every non-candidate pixel, and
$\rho=0$ recovers thresholding.
(ii)~\emph{Amplitude invariance:} if $\hat p=a_j g_j$ on $C_j$, then at $\rho=1$
the hit condition of every lesion in $C_j$ depends on the profile $g_j$ and not
on the confidence $a_j$.
(iii)~\emph{Background exclusion:} at any cutoff $c\ge t_0$ the set lies inside
$\bigcup_j C_j$.
The price is that false candidates are normalised as well. The pair
$(\rho,t_0)$ therefore sets a trade-off, and we choose it from a grid that
includes $\rho=0$.

\paragraph{Choosing the family without losing validity.}
We split the labelled pool into a tuning part and a calibration part. On the
tuning part we pick the family whose own joint threshold gives the smallest
set-to-truth ratio, then calibrate $\hat\lambda^\cup$ on the calibration part.
Conditionally on the tuning data the chosen transform is fixed. The CRC
guarantee therefore holds unchanged (Proposition~1), which would not be true if
the family were selected on the calibration data.

\section{Experiments}\label{sec:exp}

\paragraph{Protocol.}
The full protocol, including hypotheses, metrics, splits, grids and decision
rules, was committed to the repository before any result existed. Two dated
amendments, made before the corresponding results, record a change of backbone
and the addition of the simulation. Unlabelled numbers below are pre-registered;
the analysis in \S\ref{sec:posthoc} is post-hoc and labelled as such.
\emph{Data:} FracAtlas~\cite{fracatlas2023} (\FracatlasUnits\ fractured
radiographs, \FracatlasMulti\ with $\ge2$ fracture regions) and
Kvasir-SEG~\cite{kvasir2020} (\KvasirSegUnits\ polyp images, \KvasirSegMulti\
multi-polyp). Lesions are 8-connected ground-truth components labelled at native
resolution, with $\tau=0.5$. A lesion is \emph{small} if its area is below the
33rd percentile of all component areas in the dataset (a label-only definition).
\emph{Backbone:} a supervised random forest over a multiscale filter bank,
trained with 5-fold cross-fitting so that every image receives an out-of-fold
posterior. It is deliberately weak (out-of-fold Dice \FracatlasDice\ on
FracAtlas and \KvasirSegDice\ on Kvasir); see Limitations.
\emph{Splits:} \KvasirSegReps\ random tune/cal/test partitions (20/40/40\%)
per seed, with \KvasirSegSeeds\ training seeds. Baselines calibrate on tune$\cup$cal,
i.e.\ on 50\% more data than our method, which calibrates on cal only.
\emph{Metric:} the set-to-truth ratio $\sum|S|/\sum|Y|$ on test units, at a
threshold whose guarantee covers pixel FNR, lesion miss and small-lesion miss.

\paragraph{Validity.}
\ResultValidity\ (Table~\ref{tab:main}, all $\alpha$ in the supplement).
Learn-Then-Test is the most conservative method on Kvasir and at
$\alpha\le0.10$ on FracAtlas, as expected for a high-probability guarantee.

\begin{table}[t]\centering\scriptsize
\caption{Kvasir-SEG, $\alpha=0.10$, test split, mean over \KvasirSegSeeds\
seeds$\times$\KvasirSegReps\ splits. Red: a claimed risk exceeds $\alpha$
(none). $\Delta$ is the paired relative change in set size.}\label{tab:main}
\begin{tabular}{lccccr}
\toprule
Method & pixel FNR & lesion miss & small miss & set/truth & $\Delta$ vs joint-thr. \\
\midrule
Pixel CRC~\cite{angelopoulos2024crc} & 0.099 & 0.051 & 0.082 & 3.08 & -0.1\% \\
Joint, threshold & 0.098 & 0.050 & 0.081 & 3.09 & ref. \\
Joint max-form, threshold & 0.084 & 0.039 & 0.067 & 3.24 & +5.0\% \\
LTT~\cite{angelopoulos2021ltt}, threshold & 0.071 & 0.032 & 0.056 & 3.40 & +10.2\% \\
Joint, morphological~\cite{mossina2025morph} & 0.093 & 0.050 & 0.081 & 3.14 & +1.6\% \\
Joint, CCRA-S strata~\cite{luo2025ccra} & 0.092 & 0.037 & 0.066 & 3.06 & -0.8\% \\
\textbf{Ours}, joint & 0.096 & 0.049 & 0.080 & 3.11 & +0.7\% \\
\textbf{Ours}, max-form & 0.083 & 0.038 & 0.067 & 3.25 & +5.4\% \\
\bottomrule
\end{tabular}
\end{table}

\paragraph{H1: pixel guarantees leave multi-lesion units under-protected.}
At the pixel-CRC threshold, the lesion-miss rate averaged over all units stays
below $\alpha$ ($G<1$; Table~\ref{tab:gap}). On multi-lesion units, however, it
exceeds $\alpha$: on Kvasir-SEG, $G_{K\ge2}$ reaches 1.56 at $\alpha=0.10$
(verdict \KvasirSegHOne; FracAtlas \FracatlasHOne). The 95\% split intervals are
wide because only $\approx$29 multi-lesion units fall in each Kvasir calibration
split. This reproduces, on more data, the direction of the signal seen in a
smaller pilot.

\begin{table}[t]\centering\scriptsize
\caption{Gap $G=R^{\mathrm{les}}(\hat\lambda^{\mathrm{pix}})/\alpha$ on test units,
with 95\% split intervals.}\label{tab:gap}
\begin{tabular}{llccc}
\toprule
Dataset & $\alpha$ & $G$ (all) & $G\mid K{\ge}2$ & $G_{\mathrm{small}}$ \\
\midrule
FracAtlas & 0.05 & 0.64 & 0.83 [0.22, 1.64] & 0.72 [0.14, 1.46] \\
FracAtlas & 0.10 & 0.76 & 1.14 [0.59, 1.76] & 0.73 [0.33, 1.22] \\
FracAtlas & 0.15 & 0.77 & 1.26 [0.78, 1.77] & 0.81 [0.44, 1.23] \\
\midrule
Kvasir-SEG & 0.05 & 0.40 & 0.95 [0.00, 2.24] & 0.73 [0.26, 1.28] \\
Kvasir-SEG & 0.10 & 0.51 & 1.56 [0.66, 2.52] & 0.82 [0.42, 1.24] \\
Kvasir-SEG & 0.15 & 0.62 & 1.49 [0.72, 2.28] & 0.81 [0.49, 1.13] \\
\bottomrule
\end{tabular}
\end{table}

\paragraph{Joint lesion-level control is nearly free and beats existing multi-risk methods.}
At $\alpha=0.10$ the pixel-only threshold is only \FracatlasPixelCrcRel\ (FracAtlas) and \KvasirSegPixelCrcRel\ (Kvasir) smaller than the joint one, and
at most 0.7\% smaller at any $\alpha$, because on these data the pixel
constraint binds. Under the same
lesion-level guarantee, morphological sets need \KvasirSegUnionDilationRel\
(Kvasir) and \FracatlasUnionDilationRel\ (FracAtlas) larger sets at
$\alpha=0.10$. Learn-Then-Test needs \KvasirSegLttThresholdRel\ and
\FracatlasLttThresholdRel. CCRA-S strata are within $\pm2\%$ on Kvasir and 7--16\%
larger on FracAtlas.

\paragraph{H2: lesion-equalized sets on real data --- not supported.}
The pre-registered efficiency claim failed at every $\alpha$ on both datasets.
Lesion-equalized sets tied joint thresholding (Kvasir \KvasirSegOursUnionRel,
FracAtlas \FracatlasOursUnionRel\ at $\alpha=0.10$; neither smaller in $\ge95\%$
of splits). This is the behaviour the theory predicts when the lesion constraint
does not bind. There is then nothing for equalization to save, and the tuning
split often selects thresholding itself.

\paragraph{E7: when does equalization help? (controlled simulation.)}
We vary the coupling $\gamma$ between lesion size and confidence
($a\propto r^\gamma$) and the background level $b$ (Table~\ref{tab:sim}).
Lesion-equalized sets were smaller in every cell at $\gamma=0$ ($-1$ to $-4\%$),
and by up to \SimBest\ at $\gamma=\SimBestGamma$. The gain vanished at
$\gamma\ge1.5$, where the worst cell was \SimWorst. The pre-registered rule
(advantage at all $\gamma\ge1$ and increasing in $\gamma$) therefore
\SimVerdict. A post-hoc diagnostic locates the boundary. At $\gamma\ge1.5$,
76--92\% of small lesions have a peak posterior that does not exceed the image's
99th background percentile. No transform of $\hat p$ alone can then separate
them from background (Proposition~3's precondition fails), and the method falls
back to thresholding rather than doing worse.

\begin{table}[t]\centering\scriptsize
\caption{E7: change in set size (\%) of lesion-equalized vs joint thresholding,
500 splits per cell. Bold: smaller in $\ge95\%$ of splits.}\label{tab:sim}
\begin{tabular}{lccccc}
\toprule
background, $\alpha$ & $\gamma=0$ & $\gamma=0.5$ & $\gamma=1$ & $\gamma=1.5$ & $\gamma=2$ \\
\midrule
$b=0.05$, $\alpha=0.05$ & \textbf{-4.0} & \textbf{-25.7} & -19.8 & +0.7 & +0.8 \\
$b=0.05$, $\alpha=0.10$ & \textbf{-2.1} & -41.8 & \textbf{-30.7} & -0.5 & -0.5 \\
$b=0.05$, $\alpha=0.15$ & \textbf{-1.3} & \textbf{-40.0} & \textbf{-45.0} & -0.7 & +0.3 \\
$b=0.15$, $\alpha=0.05$ & \textbf{-3.6} & \textbf{-33.5} & -0.2 & +0.7 & +1.8 \\
$b=0.15$, $\alpha=0.10$ & \textbf{-2.6} & \textbf{-47.2} & -2.1 & +1.2 & +0.9 \\
$b=0.15$, $\alpha=0.15$ & \textbf{-2.0} & -38.1 & +4.0 & +1.3 & +0.8 \\
\bottomrule
\end{tabular}
\end{table}

\paragraph{Post-hoc: a conditional guarantee for multi-lesion units.}\label{sec:posthoc}
Because $L^{\mathrm{les}}$ restricted to units with $K\ge2$ is defined on a
label-determined subpopulation, it can be controlled exactly like
$L^{\mathrm{small}}$. Adding it halves the multi-lesion miss rate on Kvasir
(Table~\ref{tab:e8}), at a price of 16--28\% larger sets.

\begin{table}[t]\centering\scriptsize
\caption{Post-hoc (not pre-registered): adding the conditional risk
$L^{\mathrm{les}}\mid K\ge2$ to the joint threshold.}\label{tab:e8}
\begin{tabular}{llcccc}
\toprule
Dataset & $\alpha$ & miss$\mid K{\ge}2$, joint & miss$\mid K{\ge}2$, +cond. & set/truth, joint & set/truth, +cond. \\
\midrule
FracAtlas & 0.05 & 0.040 & 0.038 & 37.14 & 37.93 (+2\%) \\
FracAtlas & 0.10 & 0.113 & 0.092 & 27.85 & 29.89 (+7\%) \\
FracAtlas & 0.15 & 0.187 & 0.143 & 22.68 & 25.70 (+13\%) \\
Kvasir-SEG & 0.05 & 0.048 & 0.024 & 3.74 & 4.80 (+28\%) \\
Kvasir-SEG & 0.10 & 0.155 & 0.076 & 3.09 & 3.59 (+16\%) \\
Kvasir-SEG & 0.15 & 0.224 & 0.125 & 2.64 & 3.26 (+24\%) \\
\bottomrule
\end{tabular}
\end{table}

\paragraph{Limitations.}
(i)~The backbone is weak, and FracAtlas results describe a near-failing
segmenter (Dice \FracatlasDice). A GPU U-Net and the 3-D MSD Task08 pipeline are
implemented but were not run for this draft, so H2 has not been tested in the
regime where it is most plausible: a strong model whose remaining misses are
faint small lesions.
(ii)~Cross-fitted posteriors come from five models, so exchangeability is
approximate.
(iii)~Neither dataset has patient identifiers, so splits are image-wise.
(iv)~The multi-lesion subpopulation is small (29--95 units per calibration split).
(v)~The guarantees are marginal and conditional on lesion presence. None of this
is clinically validated.

\section{Conclusion}
Pixel-level conformal guarantees are carried by large lesions. On real data they
leave multi-lesion images missing lesions at up to 1.56$\alpha$. Lesion-level
risks can be controlled jointly with pixel FNR at almost no cost, and more
efficiently than with Learn-Then-Test or morphological sets. The proposed
lesion-equalized sets give large savings exactly when small lesions are present
but under-confident, and degrade gracefully to thresholding otherwise. Whether
that regime is common for strong segmenters is the open empirical question this
release is built to answer.


\bibliographystyle{splncs04}
\bibliography{refs}

\appendix
\section{Proofs}\label{app:proofs}

\noindent\textbf{Lemma 1.} If $\{S_\lambda\}$ is nested and
$L_i(\lambda)=f(S_\lambda\cap Y_i)$ with $f$ non-increasing under set inclusion, then
$\lambda\mapsto L_i(\lambda)$ is non-increasing.\\
\emph{Proof.} $\lambda\le\lambda'\Rightarrow S_\lambda\cap Y_i\subseteq S_{\lambda'}\cap Y_i$, so
$f(S_{\lambda'}\cap Y_i)\le f(S_\lambda\cap Y_i)$. Each indicator in $L^{\mathrm{les}}$ and
$L^{\mathrm{small}}$ can only flip from 1 to 0 as the covered fraction of a lesion grows. $\square$

\medskip\noindent\textbf{Lemma 2.} Let $q_k$ be the $\lceil\tau|Y^k|\rceil$-th largest score in $Y^k$.
Lesion $k$ is hit at cutoff $c$ iff $q_k\ge c$.\\
\emph{Proof.} At least $\tau|Y^k|$ pixels of $Y^k$ score $\ge c$ iff the
$\lceil\tau|Y^k|\rceil$-th largest does. $\square$

\medskip\noindent\textbf{Proposition 1 (validity with a tuned family).} Let $\mathcal T$ be a finite
set of families, and let $D_{\mathrm{tune}}$ and $D_{\mathrm{cal}}$ be disjoint, with
$D_{\mathrm{cal}}\cup\{n{+}1\}$ exchangeable conditionally on $D_{\mathrm{tune}}$. For any selection
rule $\hat T=\Psi(D_{\mathrm{tune}})$ and every risk $r$,
$\mathbb E[L^r_{n+1}(\hat\lambda^\cup)]\le\alpha$.\\
\emph{Proof.} Conditionally on $D_{\mathrm{tune}}$, $\hat T$ is a fixed label-free transform, so the
sets are nested and Lemma~1 applies. Split CRC~\cite{angelopoulos2024crc} gives
$\mathbb E[L^r_{n+1}(\hat\lambda^r)\mid D_{\mathrm{tune}}]\le\alpha$ on the subpopulation where $L^r$ is
defined. Monotonicity gives $L^r_{n+1}(\hat\lambda^\cup)\le L^r_{n+1}(\hat\lambda^r)$ pointwise. Taking
the expectation over $D_{\mathrm{tune}}$ completes the proof. $\square$

\medskip\noindent\textbf{Lemma 3.} For $\rho\in[0,1]$, $s_{\rho,t_0}\in[0,1]$. Every candidate pixel has
$s\ge t_0$, every non-candidate pixel has $s<t_0$, and $s_{0,t_0}=\hat p$.\\
\emph{Proof.} On $C_j$, $t_0\le\hat p\le P_j\le 1$ and $P_j\le P_j^\rho\le 1$. Hence
$t_0\le\hat p\le\hat p/P_j^\rho\le\hat p/P_j\le 1$. Off the candidates, $s=\hat p<t_0$. $\square$

\medskip\noindent\textbf{Proposition 2 (amplitude invariance).} If $\hat p=a_jg_j$ on $C_j$ with $a_j>0$,
then $s_{1,t_0}=g_j/\max_{C_j}g_j$ on $C_j$. Hence the hit condition of every lesion inside $C_j$
(Lemma~2) does not depend on $a_j$.\\
\emph{Proof.} $P_j=a_j\max_{C_j}g_j$; substitute into~\eqref{eq:le}. $\square$

\medskip\noindent\textbf{Proposition 3 (background exclusion).} For every $\rho$ and every cutoff
$c\ge t_0$, $\{s_{\rho,t_0}\ge c\}\subseteq\bigcup_jC_j$.\\
\emph{Proof.} Immediate from Lemma~3. $\square$

\noindent Thresholding gives no such containment. If the joint constraint needs a
thresholding cutoff $c<t_0$, every pixel in $\{c\le\hat p<t_0\}$ joins the set.
The equalized family can meet the same constraint at a cutoff $\ge t_0$ whenever
every lesion that must be hit lies inside a candidate. This is a sufficient
mechanism, not an optimality result. The loss-optimal nested family for a
lesion-level risk depends on the joint law of $(\hat p,Y)$ and remains open.

\end{document}
